\documentclass{amsart}
\usepackage{amsmath,amssymb,amsthm,hyperref}

\theoremstyle{definition}
\newtheorem{definition}{Definition}
\newtheorem{remark}{Remark}
\theoremstyle{plain}
\newtheorem{theorem}{Theorem}
\newtheorem{lemma}{Lemma}
\newtheorem{proposition}{Proposition}
\newtheorem{corollary}{Corollary}

\newcommand{\eqdef}{\mathrel{\mathop:}=}
\newcommand{\bL}{\mathbf{L}}
\newcommand{\Form}{\mathrm{Form}}
\newcommand{\fto}{\rightarrow}
\newcommand{\fbi}{\leftrightarrow}
\newcommand{\odotop}{\,\&\,}
\newcommand{\Top}{\top}
\newcommand{\Bot}{\bot}

\begin{document}

\title{A lattice-valued quantitative logic}
\author{solver-agent}
\date{}
\maketitle

\begin{abstract}
We develop a quantitative propositional logic whose truth degrees range over an
arbitrary complete residuated lattice $\bL$ (in particular over a bounded lattice
that is not linearly ordered and is not embedded in $[0,1]$).  We isolate two
complementary ``meta-degree'' constructions that replace the
integral-with-respect-to-Lebesgue-measure of the classical $[0,1]$-valued
theory:

\begin{itemize}
\item an \emph{intrinsic, $\bL$-valued} truth degree and similarity, which require
no measure at all and exist for every complete residuated lattice; and
\item a \emph{numerical}, $[0,1]$-valued truth degree obtained from a
\emph{state} (a Bosbach/Riečan averaging functional) on $\bL$, which plays exactly
the role of the uniform probability measure.
\end{itemize}

We prove for both constructions: well-definedness up to logical equivalence,
the residuated/triangle laws, that the induced object is an $\bL$-valued
similarity (resp.\ a genuine pseudometric on the set of formulas), and
compatibility with the connectives, together with a theory of divergence and
consistency degrees of theories.  We then show that the state hypothesis is sharp: an
\emph{arbitrary} order-preserving collapse $\bL\to[0,1]$ that is not a state
destroys the triangle inequality, while states (in fact top-faithful ones,
yielding genuine point-separating metrics) exist for all the standard truth-value
algebras, including non-chain ones such as the diamond $M_2$ and the product
$W_3\times W_4$.  Thus a satisfactory generalization exists for every complete
residuated lattice in the intrinsic $\bL$-valued form, and in numerical
$[0,1]$-valued form for every lattice carrying a state, with the modular/state
law being the exact additional ingredient required for the numerical metric.
\end{abstract}

\section{The algebraic setting}\label{sec:setting}

\begin{definition}[Complete residuated lattice]\label{def:rl}
A \emph{complete residuated lattice} is a structure
$\bL=(L,\wedge,\vee,\odotop,\fto,\Bot,\Top)$ such that
\begin{enumerate}
\item $(L,\wedge,\vee,\Bot,\Top)$ is a complete lattice with least element
$\Bot$ and greatest element $\Top$, with order $x\le y \iff x\wedge y = x$;
\item $(L,\odotop,\Top)$ is a commutative monoid (the \emph{multiplication}, or
\emph{strong conjunction}), with unit $\Top$;
\item $(\odotop,\fto)$ is an adjoint pair (\emph{residuation}): for all
$x,y,z\in L$,
\begin{equation}\label{eq:adj}
x\odotop y \le z \quad\Longleftrightarrow\quad x \le y\fto z .
\end{equation}
\end{enumerate}
We write $\neg x \eqdef x\fto\Bot$ and define the \emph{biresiduum}
\begin{equation}\label{eq:bi}
x\fbi y \eqdef (x\fto y)\wedge(y\fto x).
\end{equation}
\end{definition}

This class is large enough to subsume every standard truth-degree algebra used
in many-valued logic and to admit genuinely non-linear examples:

\begin{itemize}
\item the standard \L ukasiewicz, G\"odel and product algebras on $[0,1]$, and
the finite \L ukasiewicz chains $W_n$ (these recover the classical quantitative
logic of the Background paragraph);
\item every complete \emph{Heyting algebra} $(H,\wedge,\vee,\fto,\Bot,\Top)$,
where $\odotop=\wedge$ and $\fto$ is the Heyting implication
$x\fto y=\bigvee\{z: z\wedge x\le y\}$; this includes finite distributive
lattices that are \emph{not} chains, e.g.\ the diamond
$M_2=\{\Bot,a,b,\Top\}$ with $a,b$ incomparable;
\item any direct product $\bL_1\times\bL_2$ of residuated lattices, with
componentwise operations; e.g.\ $W_3\times W_4$, a $12$-element residuated lattice
that is not a chain.
\end{itemize}

We collect the elementary residuation facts used throughout.  Each is a standard
consequence of \eqref{eq:adj}; we prove the ones we need from scratch.

\begin{lemma}[Residuation calculus]\label{lem:res}
For all $x,y,z\in L$:
\begin{enumerate}
\item\label{it:r-mono} $\odotop$ is order-preserving in each argument;
$\fto$ is order-preserving in the second and order-reversing in the first.
\item\label{it:r-unit} $x\fto x=\Top$ and $\Top\fto x=x$.
\item\label{it:r-mp} $x\odotop(x\fto y)\le y$ (modus ponens).
\item\label{it:r-trans} $(x\fto y)\odotop(y\fto z)\le x\fto z$.
\item\label{it:bi-trans} $(x\fbi y)\odotop(y\fbi z)\le x\fbi z$.
\item\label{it:bi-refl} $x\fbi x=\Top$, \ $x\fbi y=y\fbi x$, and $x\fbi y=\Top$
iff $x=y$.
\end{enumerate}
\end{lemma}

\begin{proof}
\eqref{it:r-mono} If $a\le b$ then $a\odotop c\le b\odotop c$: indeed
$b\odotop c\le b\odotop c$ gives $b\le c\fto(b\odotop c)$ by \eqref{eq:adj}, and
$a\le b\le c\fto(b\odotop c)$ gives $a\odotop c\le b\odotop c$ again by
\eqref{eq:adj}.  Monotonicity of $\fto$ in the second argument: if $y\le y'$ and
$z=x\fto y$ then $z\odotop x\le y\le y'$, so $z\le x\fto y'$.  Antitonicity in the
first: if $x\le x'$ and $z=x'\fto y$ then $z\odotop x\le z\odotop x'\le y$, hence
$z\le x\fto y$.

\eqref{it:r-unit} From $x\odotop\Top=x\le x$ and \eqref{eq:adj}, $\Top\le x\fto x$,
so $x\fto x=\Top$.  And $\Top\fto x$: from \eqref{eq:adj},
$z\le\Top\fto x \iff z=z\odotop\Top\le x$, so $\Top\fto x=x$.

\eqref{it:r-mp} $x\fto y\le x\fto y$ gives, by \eqref{eq:adj},
$(x\fto y)\odotop x\le y$, i.e.\ $x\odotop(x\fto y)\le y$ by commutativity.

\eqref{it:r-trans} By \eqref{it:r-mp} (twice) and \eqref{it:r-mono},
\[
\bigl[(x\fto y)\odotop(y\fto z)\bigr]\odotop x
=(y\fto z)\odotop\bigl[(x\fto y)\odotop x\bigr]
\le (y\fto z)\odotop y \le z .
\]
By \eqref{eq:adj} this yields $(x\fto y)\odotop(y\fto z)\le x\fto z$.

\eqref{it:bi-trans} Using \eqref{it:r-mono}, $\eqref{eq:bi}$ and \eqref{it:r-trans},
\begin{align*}
(x\fbi y)\odotop(y\fbi z)
&\le (x\fto y)\odotop(y\fto z)\le x\fto z,\\
(x\fbi y)\odotop(y\fbi z)
&\le (y\fto x)\odotop(z\fto y)=(z\fto y)\odotop(y\fto x)\le z\fto x.
\end{align*}
Hence $(x\fbi y)\odotop(y\fbi z)\le (x\fto z)\wedge(z\fto x)=x\fbi z$.

\eqref{it:bi-refl} $x\fbi x=(x\fto x)\wedge(x\fto x)=\Top$ by \eqref{it:r-unit};
symmetry is immediate from \eqref{eq:bi}.  If $x\fbi y=\Top$ then
$x\fto y=\Top$ and $y\fto x=\Top$; from $x\fto y=\Top$ and \eqref{eq:adj}
(with the left side $\Top\odotop x=x$) we get $x\le y$, and symmetrically
$y\le x$, so $x=y$.  The converse is $x\fbi x=\Top$.
\end{proof}

\section{Formulas and their induced truth functions}\label{sec:formulas}

Fix a countable set of propositional variables $\{p_1,p_2,\dots\}$.

\begin{definition}[Formulas]\label{def:form}
The set $\Form$ of \emph{formulas} is the least set containing all variables and
the constants $\overline\Bot,\overline\Top$, and closed under the binary
connectives $\wedge,\vee,\odotop$ (multiplicative conjunction) and $\fto$.  We use
$\neg\varphi\eqdef\varphi\fto\overline\Bot$ and
$\varphi\fbi\psi\eqdef(\varphi\fto\psi)\wedge(\psi\fto\varphi)$ as abbreviations.
For $\varphi\in\Form$ let $\mathrm{Var}(\varphi)$ be its (finite) set of variables;
write $\Form_n$ for the formulas with $\mathrm{Var}(\varphi)\subseteq\{p_1,\dots,p_n\}$.
\end{definition}

\begin{definition}[Valuations and induced functions]\label{def:val}
A \emph{valuation} is a map $v:\{p_i\}\to L$.  It extends uniquely to
$\overline v:\Form\to L$ by interpreting each connective by the corresponding
operation of $\bL$ and $\overline\Bot,\overline\Top$ by $\Bot,\Top$.  For
$\varphi\in\Form_n$ the \emph{induced truth function}
$\overline\varphi:L^n\to L$ is
$\overline\varphi(a_1,\dots,a_n)\eqdef\overline v(\varphi)$ for any valuation $v$
with $v(p_i)=a_i$ ($1\le i\le n$); this is well defined since
$\overline v(\varphi)$ depends only on $v\restriction\mathrm{Var}(\varphi)$.
\end{definition}

\begin{definition}[Semantic equivalence]\label{def:eqv}
Formulas $\varphi,\psi$ are \emph{$\bL$-equivalent}, written
$\varphi\approx\psi$, if $\overline\varphi=\overline\psi$ as functions on
$L^{m}$, where $m$ is large enough that
$\mathrm{Var}(\varphi)\cup\mathrm{Var}(\psi)\subseteq\{p_1,\dots,p_m\}$.
Equivalently (by Lemma~\ref{lem:res}\eqref{it:bi-refl}),
$\varphi\approx\psi \iff \overline{\varphi\fbi\psi}\equiv\Top$, i.e.\ the induced
function of $\varphi\fbi\psi$ is constantly $\Top$.  Plainly $\approx$ is an
equivalence relation on $\Form$.
\end{definition}

Throughout, when comparing functions of finitely many variables we always pad to
a common variable set $\{p_1,\dots,p_m\}$; adding dummy variables does not change
any quantity below, because all our aggregations are over product structures and
the marginal in a dummy variable is constant.

\section{Two notions of meta-degree}\label{sec:meta}

The classical theory averages $\overline\varphi:[0,1]^n\to[0,1]$ against Lebesgue
measure on $[0,1]^n$.  Two structural facts make this work: \emph{(i)} $[0,1]$
carries a translation that lets one form an average; \emph{(ii)} that average is
\emph{monotone} and \emph{modular} (additive).  A general lattice has neither a
canonical addition nor a canonical measure.  We therefore separate the two roles.

\subsection{The intrinsic ($\bL$-valued) meta-degree}

For non-linear $\bL$ there is a canonical aggregation that needs no measure: the
infimum and supremum.  These give a meta-degree and similarity valued in $\bL$
itself.

\begin{definition}[Intrinsic meta-degree and similarity]\label{def:intrinsic}
For $\varphi\in\Form_n$ put
\begin{equation}\label{eq:intr}
\tau_\wedge(\varphi)\eqdef\bigwedge_{a\in L^n}\overline\varphi(a)\in L,
\qquad
\tau_\vee(\varphi)\eqdef\bigvee_{a\in L^n}\overline\varphi(a)\in L .
\end{equation}
For $\varphi,\psi\in\Form$ define the \emph{intrinsic similarity}
\begin{equation}\label{eq:Esim}
E(\varphi,\psi)\eqdef
\bigwedge_{a\in L^{m}}\overline{(\varphi\fbi\psi)}(a)\;=\;\tau_\wedge(\varphi\fbi\psi)\in L,
\end{equation}
with $m$ a common variable bound.  (Infima/suprema exist because $\bL$ is
complete.)  We call $E(\varphi,\psi)$ the \emph{degree to which $\varphi$ and
$\psi$ are interchangeable}; it lives in $\bL$.
\end{definition}

\begin{theorem}[Intrinsic similarity]\label{thm:intrinsic}
The map $E:\Form\times\Form\to L$ of \eqref{eq:Esim} is well defined,
independent of the chosen common variable bound $m$, and an $\bL$-valued
similarity:
\begin{enumerate}
\item\label{it:E-refl} $E(\varphi,\varphi)=\Top$;
\item\label{it:E-sym} $E(\varphi,\psi)=E(\psi,\varphi)$;
\item\label{it:E-trans} $E(\varphi,\psi)\odotop E(\psi,\chi)\le E(\varphi,\chi)$
(residuated transitivity);
\item\label{it:E-sep} $E(\varphi,\psi)=\Top$ iff $\varphi\approx\psi$;
\item\label{it:E-cong} $E$ is a congruence for the connectives: for
$*\in\{\wedge,\vee,\odotop,\fto\}$,
\[
E(\varphi_1,\psi_1)\odotop E(\varphi_2,\psi_2)\ \le\
E(\varphi_1*\varphi_2,\ \psi_1*\psi_2).
\]
\end{enumerate}
In particular $E$ descends to $\Form/{\approx}$ and there
$E(\varphi,\psi)=\Top\iff\varphi=\psi$.
\end{theorem}

\begin{proof}
Independence of $m$: adding a dummy variable $p_{m+1}$ replaces $L^m$ by
$L^{m+1}=L^m\times L$, and $\overline{(\varphi\fbi\psi)}$ does not depend on the
last coordinate, so the infimum is unchanged.

\eqref{it:E-refl}, \eqref{it:E-sym} follow from
Lemma~\ref{lem:res}\eqref{it:bi-refl} applied pointwise:
$\overline{(\varphi\fbi\varphi)}(a)=\overline\varphi(a)\fbi\overline\varphi(a)=\Top$
for all $a$, and $\overline{(\varphi\fbi\psi)}=\overline{(\psi\fbi\varphi)}$.

\eqref{it:E-trans} For each $a\in L^m$,
Lemma~\ref{lem:res}\eqref{it:bi-trans} gives
\[
\overline{(\varphi\fbi\psi)}(a)\odotop\overline{(\psi\fbi\chi)}(a)
=\bigl(\overline\varphi(a)\fbi\overline\psi(a)\bigr)\odotop
\bigl(\overline\psi(a)\fbi\overline\chi(a)\bigr)
\le \overline\varphi(a)\fbi\overline\chi(a)
=\overline{(\varphi\fbi\chi)}(a).
\]
Now $E(\varphi,\psi)\le\overline{(\varphi\fbi\psi)}(a)$ and
$E(\psi,\chi)\le\overline{(\psi\fbi\chi)}(a)$ for every $a$, so by
Lemma~\ref{lem:res}\eqref{it:r-mono},
$E(\varphi,\psi)\odotop E(\psi,\chi)\le\overline{(\varphi\fbi\chi)}(a)$ for every
$a$; taking the infimum over $a$ gives \eqref{it:E-trans}.

\eqref{it:E-sep} $E(\varphi,\psi)=\Top$ iff
$\overline{(\varphi\fbi\psi)}(a)=\Top$ for all $a$ (an infimum equals $\Top$ iff
every term is $\Top$), iff $\overline\varphi(a)=\overline\psi(a)$ for all $a$
(Lemma~\ref{lem:res}\eqref{it:bi-refl}), iff $\varphi\approx\psi$.

\eqref{it:E-cong} We first prove the pointwise statement for elements of $L$:
for all $x_i,y_i\in L$ and $*\in\{\wedge,\vee,\odotop,\fto\}$,
\begin{equation}\label{eq:point-cong}
(x_1\fbi y_1)\odotop(x_2\fbi y_2)\ \le\ (x_1*x_2)\fbi(y_1*y_2).
\end{equation}
\emph{Case $*=\odotop$.}  Put $e_i=x_i\fbi y_i$.  By
Lemma~\ref{lem:res}\eqref{it:r-mp}, $e_i\odotop x_i\le e_i\odotop x_i$ and
$x_i\fto y_i\ge e_i$ give, using
$x_i\odotop(x_i\fto y_i)\le y_i$ and monotonicity,
$x_i\odotop e_i\le y_i$; symmetrically $y_i\odotop e_i\le x_i$.  Then
\[
(x_1\odotop x_2)\odotop(e_1\odotop e_2)
=(x_1\odotop e_1)\odotop(x_2\odotop e_2)\le y_1\odotop y_2,
\]
so by \eqref{eq:adj} $e_1\odotop e_2\le (x_1\odotop x_2)\fto(y_1\odotop y_2)$;
symmetrically $e_1\odotop e_2\le(y_1\odotop y_2)\fto(x_1\odotop x_2)$, proving
\eqref{eq:point-cong} for $\odotop$.
\emph{Case $*=\fto$.}  We claim
$e_1\odotop e_2\le(x_1\fto x_2)\fto(y_1\fto y_2)$.  By \eqref{eq:adj} it suffices
that $e_1\odotop e_2\odotop(x_1\fto x_2)\le y_1\fto y_2$, i.e.\ by \eqref{eq:adj}
again that $e_1\odotop e_2\odotop(x_1\fto x_2)\odotop y_1\le y_2$.  Now
$y_1\odotop e_1\le x_1$ (shown above), and $x_1\odotop(x_1\fto x_2)\le x_2$
(Lemma~\ref{lem:res}\eqref{it:r-mp}), and $x_2\odotop e_2\le y_2$, so using
commutativity and monotonicity,
\[
e_1\odotop e_2\odotop(x_1\fto x_2)\odotop y_1
\le e_2\odotop(x_1\fto x_2)\odotop x_1
\le e_2\odotop x_2\le y_2 .
\]
By symmetry $e_1\odotop e_2\le(y_1\fto y_2)\fto(x_1\fto x_2)$, giving
\eqref{eq:point-cong} for $\fto$.
\emph{Cases $*=\wedge,\vee$.}  Since $e_i\le x_i\fto y_i$ and
$e_i\le y_i\fto x_i$, by modus ponens $x_i\odotop e_i\le y_i$ and
$y_i\odotop e_i\le x_i$.  Put $f=e_1\odotop e_2\le e_1,e_2$.  Then
$x_i\odotop f\le x_i\odotop e_i\le y_i$ for $i=1,2$, whence
$(x_1\wedge x_2)\odotop f\le x_i\odotop f\le y_i$, so
$(x_1\wedge x_2)\odotop f\le y_1\wedge y_2$ and by \eqref{eq:adj}
$f\le (x_1\wedge x_2)\fto(y_1\wedge y_2)$; the symmetric bound and the
$\vee$-case ($(x_1\vee x_2)\odotop f=(x_1\odotop f)\vee(x_2\odotop f)$, using that
$\odotop$ distributes over $\vee$ which follows from \eqref{eq:adj} and
completeness) are identical.  This proves \eqref{eq:point-cong}.

Finally evaluate \eqref{eq:point-cong} at $x_i=\overline{\varphi_i}(a)$,
$y_i=\overline{\psi_i}(a)$ and bound each $x_i\fbi y_i$ below by
$E(\varphi_i,\psi_i)$; monotonicity of $\odotop$ gives
$E(\varphi_1,\psi_1)\odotop E(\varphi_2,\psi_2)\le
\overline{(\varphi_1*\varphi_2)}(a)\fbi\overline{(\psi_1*\psi_2)}(a)$ for all
$a$, and taking the infimum gives \eqref{it:E-cong}.
\end{proof}

\begin{remark}[Codomain of the intrinsic degree]\label{rem:intr-codomain}
Theorem~\ref{thm:intrinsic} answers part (b) in its first form: the
\emph{intrinsic} meta-degree of interchangeability takes values in the lattice
$\bL$ itself; no measure is needed.  Its ``metric'' is the residuated triangle law
\eqref{it:E-trans}, the exact lattice analogue of the triangle inequality (it
becomes the usual triangle inequality after applying $-\log$ to a real residuated
lattice).  This is the largest invariant available with no extra structure on
$\bL$, and it is non-trivial precisely because $\bL$ may be non-linear.
\end{remark}

\subsection{States: the role of the uniform measure}

To obtain a \emph{numerical} truth degree in $[0,1]$ as in the classical theory,
we replace ``Lebesgue measure followed by integration'' by a single averaging
functional, a \emph{state}.

\begin{definition}[State]\label{def:state}
A \emph{state} on $\bL$ is a map $s:L\to[0,1]$ with
\begin{enumerate}
\item[(S1)] $s(\Bot)=0$, $s(\Top)=1$;
\item[(S2)] $s$ is order-preserving: $x\le y\Rightarrow s(x)\le s(y)$;
\item[(S3)] (modular/valuation law) $s(x\vee y)+s(x\wedge y)=s(x)+s(y)$ for all $x,y$;
\item[(S4)] ($\&$-subadditivity) $s(x)+s(y)\le 1+s(x\odotop y)$ for all $x,y$.
\end{enumerate}
\end{definition}

\begin{remark}[Existence and the classical case]\label{rem:state-exist}
Conditions (S1)--(S4) hold for the canonical averaging functional of the
classical theory.  On $[0,1]$ with the \L ukasiewicz operations
$x\odotop y=\max(0,x+y-1)$, take $s=\mathrm{id}$: (S1),(S2) are clear,
(S3) is $\max(x,y)+\min(x,y)=x+y$, and (S4) is $x+y\le 1+\max(0,x+y-1)$, true
since $\max(0,x+y-1)\ge x+y-1$.  For any finite residuated chain $W_k$ the same
$s(x)=x$ is a state.  States also exist for all non-chain examples: by
Proposition~\ref{prop:exist} every finite Heyting algebra, every finite
MV-algebra, and their products carry the height-function state.  Concretely, on
the diamond $M_2$ the height state $s(\Bot)=0$, $s(a)=s(b)=\tfrac12$,
$s(\Top)=1$ satisfies (S1)--(S4): (S3) is
$s(a\vee b)+s(a\wedge b)=s(\Top)+s(\Bot)=1=s(a)+s(b)$, and (S4) with
$\odotop=\wedge$ reads $s(x)+s(y)\le 1+s(x\wedge y)$, which by (S3) is
$s(x\vee y)\le1$.  Existence and its sharp boundary are treated in
\S\ref{sec:obstr}; the intrinsic theory of the previous subsection requires no
state at all.
\end{remark}

We now choose the averaging measure on valuations.  Because $\overline\varphi$
factors through finitely many coordinates, the natural ``uniform'' object is a
sequence of probability spaces.

\begin{definition}[Averaged numerical meta-degree]\label{def:tau-state}
Fix a state $s$ on $\bL$ and, for each finite $L_0\subseteq L$ that generates
the values needed below, the uniform counting measure.  Concretely, when $\bL$ is
finite let $\mu_n$ be the uniform probability measure on $L^n$
($\mu_n(\{a\})=|L|^{-n}$); when $\bL$ is infinite assume a fixed Borel probability
measure $\mu_n$ on $L^n$ that is a product $\mu_1^{\otimes n}$ of a fixed measure
$\mu_1$ on $L$ for which $a\mapsto s(\overline\varphi(a))$ is $\mu_n$-integrable
for every $\varphi$ (e.g.\ Lebesgue/uniform when $L=[0,1]$).  Define the
\emph{numerical truth degree}
\begin{equation}\label{eq:tau-s}
\tau_s(\varphi)\eqdef \int_{L^n} s\bigl(\overline\varphi(a)\bigr)\,d\mu_n(a)\in[0,1],
\qquad \varphi\in\Form_n,
\end{equation}
and the \emph{logic distance}
\begin{equation}\label{eq:rho-s}
\rho_s(\varphi,\psi)\eqdef 1-\tau_s(\varphi\fbi\psi)
=\int_{L^m}\bigl[\,1-s\bigl(\overline{(\varphi\fbi\psi)}(a)\bigr)\bigr]\,d\mu_m(a).
\end{equation}
\end{definition}

The integrand $a\mapsto s(\overline\varphi(a))$ is $\mu_n$-measurable: in the
finite case it is a finite-valued function; in the assumed infinite case it is
measurable because $s$ is monotone (hence Borel) and $\overline\varphi$ is a
composition of the (assumed Borel) operations of $\bL$.

\begin{lemma}[Well-definedness and basic shape of $\tau_s$]\label{lem:tau-wd}
With $\mu_n=\mu_1^{\otimes n}$ as above:
\begin{enumerate}
\item\label{it:tau-dummy} $\tau_s(\varphi)$ is independent of the number of dummy
variables: padding $\varphi\in\Form_n$ to $\Form_{n+1}$ leaves
\eqref{eq:tau-s} unchanged.
\item\label{it:tau-eqv} $\varphi\approx\psi\Rightarrow\tau_s(\varphi)=\tau_s(\psi)$
and $\rho_s(\varphi,\psi)=0$.
\item\label{it:tau-range} $0\le\tau_s(\varphi)\le1$, with
$\tau_s(\overline\Top)=1$, $\tau_s(\overline\Bot)=0$.
\end{enumerate}
\end{lemma}

\begin{proof}
\eqref{it:tau-dummy} Since $\overline\varphi$ ignores $p_{n+1}$, Fubini for the
product measure gives
$\int_{L^{n+1}}s(\overline\varphi)\,d\mu_{n+1}
=\int_{L^n}s(\overline\varphi)\,d\mu_n\cdot\mu_1(L)
=\int_{L^n}s(\overline\varphi)\,d\mu_n$.
\eqref{it:tau-eqv} If $\varphi\approx\psi$ then $\overline\varphi=\overline\psi$
pointwise, so the integrands agree; and $\overline{(\varphi\fbi\psi)}\equiv\Top$
gives $s(\overline{(\varphi\fbi\psi)})\equiv1$, so
$\tau_s(\varphi\fbi\psi)=1$ and $\rho_s(\varphi,\psi)=0$.
\eqref{it:tau-range} Immediate from $0\le s\le1$, $s(\Top)=1$, $s(\Bot)=0$ and
$\mu_n(L^n)=1$.
\end{proof}

\section{The induced logic pseudometric}\label{sec:metric}

\begin{theorem}[$\rho_s$ is a logic pseudometric]\label{thm:pseudometric}
Let $s$ be a state on $\bL$.  Then $\rho_s$ of \eqref{eq:rho-s} is a pseudometric
on $\Form$:
\begin{enumerate}
\item\label{it:rho-zero} $\rho_s(\varphi,\varphi)=0$;
\item\label{it:rho-sym} $\rho_s(\varphi,\psi)=\rho_s(\psi,\varphi)$;
\item\label{it:rho-tri} $\rho_s(\varphi,\chi)\le\rho_s(\varphi,\psi)+\rho_s(\psi,\chi)$.
\end{enumerate}
Moreover $\rho_s(\varphi,\psi)=0$ whenever $\varphi\approx\psi$, so $\rho_s$
descends to a pseudometric on $\Form/{\approx}$.
\end{theorem}

\begin{proof}
\eqref{it:rho-zero} By Lemma~\ref{lem:res}\eqref{it:bi-refl},
$\overline{(\varphi\fbi\varphi)}\equiv\Top$, so the integrand in \eqref{eq:rho-s}
is identically $0$.
\eqref{it:rho-sym} $\overline{(\varphi\fbi\psi)}=\overline{(\psi\fbi\varphi)}$
pointwise (Lemma~\ref{lem:res}\eqref{it:bi-refl}), so the integrands agree.
\eqref{it:rho-tri} We prove the \emph{pointwise} triangle inequality
\begin{equation}\label{eq:point-tri}
1-s\bigl(\overline\varphi(a)\fbi\overline\chi(a)\bigr)
\ \le\
\bigl[1-s(\overline\varphi(a)\fbi\overline\psi(a))\bigr]
+\bigl[1-s(\overline\psi(a)\fbi\overline\chi(a))\bigr]
\end{equation}
for every $a\in L^m$, and then integrate.  Write $x=\overline\varphi(a)$,
$y=\overline\psi(a)$, $z=\overline\chi(a)$, and $u=x\fbi y$, $w=y\fbi z$.  By
Lemma~\ref{lem:res}\eqref{it:bi-trans}, $u\odotop w\le x\fbi z$, so by (S2),
\begin{equation}\label{eq:tri-a}
s(u\odotop w)\le s(x\fbi z).
\end{equation}
By (S4) applied to $u,w$,
\begin{equation}\label{eq:tri-b}
s(u)+s(w)\le 1+s(u\odotop w).
\end{equation}
Combining \eqref{eq:tri-b} and \eqref{eq:tri-a}:
$s(u)+s(w)\le 1+s(x\fbi z)$, i.e.
\[
s(x\fbi y)+s(y\fbi z)\le 1+s(x\fbi z),
\]
which rearranges to exactly \eqref{eq:point-tri}.  Integrating \eqref{eq:point-tri}
over $a\in L^m$ against the probability measure $\mu_m$ (and using linearity of the
integral and $\mu_m(L^m)=1$) gives \eqref{it:rho-tri}.

The last clause is Lemma~\ref{lem:tau-wd}\eqref{it:tau-eqv}.
\end{proof}

\begin{remark}[Why a state, not an arbitrary collapse]\label{rem:why-state}
Both (S3)/(S4) and (S2) are used in \eqref{eq:tri-a}--\eqref{eq:tri-b}.  Dropping
the valuation/subadditivity axioms genuinely breaks \eqref{it:rho-tri}; see
\S\ref{sec:obstr}.  Thus the state law is not a convenience but the precise
hypothesis under which the integral construction of the $[0,1]$-valued theory
survives the passage to a non-linear $\bL$.
\end{remark}

\begin{proposition}[Compatibility with connectives]\label{prop:lipschitz}
Let $s$ be a state on $\bL$.  For all $\varphi_1,\varphi_2,\psi_1,\psi_2\in\Form$
and $*\in\{\wedge,\vee,\odotop,\fto\}$,
\begin{equation}\label{eq:lip}
\rho_s(\varphi_1*\varphi_2,\ \psi_1*\psi_2)\ \le\
\rho_s(\varphi_1,\psi_1)+\rho_s(\varphi_2,\psi_2).
\end{equation}
In particular each connective is non-expansive (uniformly continuous) for
$\rho_s$, and $\rho_s(\neg\varphi,\neg\psi)\le\rho_s(\varphi,\psi)$.
\end{proposition}

\begin{proof}
Fix $a$ and write $x_i=\overline{\varphi_i}(a)$, $y_i=\overline{\psi_i}(a)$, and
$e_i=x_i\fbi y_i$.  By the pointwise congruence \eqref{eq:point-cong} of
Theorem~\ref{thm:intrinsic}\eqref{it:E-cong},
$e_1\odotop e_2\le(x_1*x_2)\fbi(y_1*y_2)$, so by (S2),
\[
s\bigl((x_1*x_2)\fbi(y_1*y_2)\bigr)\ \ge\ s(e_1\odotop e_2).
\]
By (S4), $s(e_1\odotop e_2)\ge s(e_1)+s(e_2)-1$.  Hence
\[
1-s\bigl((x_1*x_2)\fbi(y_1*y_2)\bigr)
\le 1-(s(e_1)+s(e_2)-1)
=\bigl[1-s(e_1)\bigr]+\bigl[1-s(e_2)\bigr].
\]
Integrating over $a$ against $\mu_m$ gives \eqref{eq:lip}.  For negation take
$*=\fto$, $\varphi_2=\psi_2=\overline\Bot$ (so $\rho_s$-term $=0$): then
$\varphi_1\fto\overline\Bot=\neg\varphi_1$ and we get
$\rho_s(\neg\varphi_1,\neg\psi_1)\le\rho_s(\varphi_1,\psi_1)$.
\end{proof}

\section{Divergence and consistency degrees of theories}\label{sec:theories}

Let $\Gamma\subseteq\Form$ be a \emph{theory} (a set of formulas).

\begin{definition}[Truth, divergence, consistency]\label{def:cons}
Fix a state $s$.  The \emph{degree of truth of $\Gamma$} and the
\emph{divergence degree} of $\Gamma$ are
\begin{equation}\label{eq:cons}
\tau_s(\Gamma)\eqdef \inf_{\varphi\in\Gamma}\tau_s(\varphi)\in[0,1],
\qquad
\mathrm{div}_s(\Gamma)\eqdef\sup_{\varphi,\psi\in\Gamma}\rho_s(\varphi,\psi)\in[0,1].
\end{equation}
We say $\Gamma$ is \emph{$\rho_s$-consistent} if $\mathrm{div}_s(\Gamma)<1$, and
\emph{fully inconsistent} if $\mathrm{div}_s(\Gamma)=1$.  The
\emph{consistency degree} is $\mathrm{cons}_s(\Gamma)\eqdef 1-\mathrm{div}_s(\Gamma)$.
\end{definition}

\begin{proposition}[Properties of the divergence/consistency degree]\label{prop:cons}
For all theories $\Gamma,\Gamma'$ and state $s$:
\begin{enumerate}
\item\label{it:c-mono} $\Gamma\subseteq\Gamma'\Rightarrow
\mathrm{div}_s(\Gamma)\le\mathrm{div}_s(\Gamma')$ and
$\mathrm{cons}_s(\Gamma)\ge\mathrm{cons}_s(\Gamma')$.
\item\label{it:c-inv} $\mathrm{div}_s$ and $\mathrm{cons}_s$ depend on $\Gamma$ only
through $\{[\varphi]_\approx:\varphi\in\Gamma\}\subseteq\Form/{\approx}$.
\item\label{it:c-tri} For any $\theta\in\Form$,
$\mathrm{div}_s(\Gamma)\le 2\sup_{\varphi\in\Gamma}\rho_s(\varphi,\theta)$;
hence if some $\theta$ is $\rho_s$-close to all of $\Gamma$ then $\Gamma$ is
consistent.  In particular every theory all of whose members are
$\bL$-equivalent has divergence $0$.
\item\label{it:c-singleton} A single formula $\Gamma=\{\varphi\}$ has
$\mathrm{div}_s(\Gamma)=0$, i.e.\ is maximally consistent.
\item\label{it:c-full} $\Gamma$ is fully inconsistent iff
$\sup_{\varphi,\psi\in\Gamma}\rho_s(\varphi,\psi)=1$; e.g.\ if $\Gamma$ contains
two formulas $\varphi,\psi$ with $\overline{(\varphi\fbi\psi)}\equiv\Bot$ and
$s(\Bot)=0$, then $\rho_s(\varphi,\psi)=1$ and $\Gamma$ is fully inconsistent.
\end{enumerate}
\end{proposition}

\begin{proof}
\eqref{it:c-mono} Supremum over a larger index set.
\eqref{it:c-inv} By Lemma~\ref{lem:tau-wd}\eqref{it:tau-eqv}, $\rho_s$ and
$\tau_s$ are constant on $\approx$-classes.
\eqref{it:c-tri} For $\varphi,\psi\in\Gamma$, the triangle inequality
(Theorem~\ref{thm:pseudometric}\eqref{it:rho-tri}) gives
$\rho_s(\varphi,\psi)\le\rho_s(\varphi,\theta)+\rho_s(\theta,\psi)\le
2\sup_{\xi\in\Gamma}\rho_s(\xi,\theta)$; take the supremum over $\varphi,\psi$.
If all members are $\approx$-equivalent, pick $\theta\in\Gamma$; then each
$\rho_s(\varphi,\theta)=0$ and $\mathrm{div}_s(\Gamma)=0$.
\eqref{it:c-singleton} $\rho_s(\varphi,\varphi)=0$
(Theorem~\ref{thm:pseudometric}\eqref{it:rho-zero}).
\eqref{it:c-full} The first statement is the definition; for the example,
$\overline{(\varphi\fbi\psi)}\equiv\Bot$ gives integrand $1-s(\Bot)=1$, so
$\rho_s(\varphi,\psi)=1$.
\end{proof}

\begin{remark}[Approximate reasoning]\label{rem:approx}
Proposition~\ref{prop:lipschitz} is the engine of approximate reasoning: if a
derivation uses connectives finitely many times and the hypotheses are perturbed
by total $\rho_s$-error $\varepsilon$, the conclusion is perturbed by at most a
constant times $\varepsilon$ (the constant being the number of connective
applications), exactly as in the classical $[0,1]$-valued logic.  All of this is
now available over a non-linear truth-value lattice.
\end{remark}

\section{Faithfulness, sharpness, and obstructions}\label{sec:obstr}

We now delimit precisely what generalizes.  Three facts emerge: the
modular/state law cannot be weakened to mere monotonicity; states (hence
pseudometrics) exist for all the standard truth-value algebras including
non-chain ones; and a mild faithfulness condition on the state turns $\rho_s$
into a genuine point-separating metric.

\begin{proposition}[An arbitrary monotone collapse breaks the triangle inequality]
\label{prop:obstr-mono}
There is a finite Heyting algebra $\bL$ and a monotone map $v:L\to[0,1]$ with
$v(\Bot)=0$, $v(\Top)=1$ \emph{(but $v$ not a state)} such that the function
$\rho_v(\varphi,\psi)=1-v(\overline{(\varphi\fbi\psi)})$, evaluated at a single
valuation, violates the triangle inequality.  Hence the modular law (S3) cannot
be dropped: monotonicity alone does not yield a pseudometric.
\end{proposition}

\begin{proof}
Let $\bL=M_2=\{\Bot,a,b,\Top\}$ be the diamond Heyting algebra
($\odotop=\wedge$, $a,b$ incomparable, $a\wedge b=\Bot$, $a\vee b=\Top$).  Its
biresiduum, computed from the Heyting implication
$x\fto y=\bigvee\{z:z\wedge x\le y\}$, gives
\[
a\fbi b=\Bot,\qquad a\fbi\Top=a,\qquad b\fbi\Top=b .
\]
Take the monotone $v$ with $v(\Bot)=0$, $v(a)=v(b)=\tfrac23$, $v(\Top)=1$.  With
$(x,y,z)=(a,\Top,b)$,
\[
1-v(a\fbi b)=1-v(\Bot)=1,\qquad
\bigl[1-v(a\fbi\Top)\bigr]+\bigl[1-v(\Top\fbi b)\bigr]
=(1-\tfrac23)+(1-\tfrac23)=\tfrac23,
\]
and $1>\tfrac23$, so the pointwise triangle inequality fails.  This $v$ violates
(S3): $v(a\vee b)+v(a\wedge b)=v(\Top)+v(\Bot)=1$ while $v(a)+v(b)=\tfrac43$.  By
contrast the modular value $v(a)=v(b)=\tfrac12$ is a state and, by
Theorem~\ref{thm:pseudometric}, gives a pseudometric; the failure above is
attributable exactly to the loss of (S3).
\end{proof}

\begin{proposition}[States exist for the standard algebras]\label{prop:exist}
Let $\bL$ be either a finite Heyting algebra \emph{(}so $\odotop=\wedge$, the
underlying lattice finite distributive, possibly non-chain\emph{)} or a finite
MV-algebra, or a finite direct product of such.  Let $h$ be the height function
$h(x)=$ the common length of every maximal chain from $\Bot$ to $x$ \emph{(}well
defined because each such $\bL$ is graded\emph{)}, and set
$s(x)\eqdef h(x)/h(\Top)$.  Then:
\begin{enumerate}
\item\label{it:exist-state} $s$ is a state on $\bL$, hence $\rho_s$ is a
pseudometric \emph{(Theorem~\ref{thm:pseudometric})};
\item\label{it:exist-topf} $s$ is \emph{top-faithful}: $s(e)=1\iff e=\Top$.
\end{enumerate}
In particular states, hence numerical logic pseudometrics, exist for all the
non-chain examples such as $M_2$ and $W_3\times W_4$.
\end{proposition}

\begin{proof}
\emph{Gradedness and the rank identity.}  Every finite distributive lattice is
modular, every chain is modular, and a finite product of modular graded lattices
is modular and graded with height the coordinatewise sum of factor heights.  In a
finite modular lattice every maximal chain between two comparable elements has the
same length (Jordan--Dedekind condition), so $h$ is well defined with
$h(\Bot)=0$, $h(y)=h(x)+1$ whenever $y$ covers $x$, and $h$ satisfies the modular
rank identity $h(x\vee y)+h(x\wedge y)=h(x)+h(y)$ for all $x,y$.  Dividing by
$h(\Top)>0$ gives $s(\Bot)=0$, $s(\Top)=1$ (S1), monotonicity (S2) from
$x\le y\Rightarrow h(x)\le h(y)$, and the modular law (S3).

\emph{(S4).}  For a Heyting algebra $\odotop=\wedge$, so (S4) reads
$s(x)+s(y)\le 1+s(x\wedge y)$; by (S3) this is $s(x\vee y)\le 1$, which holds.
For an MV-algebra the rank $s$ satisfies the MV-state identity
$s(x)+s(y)=s(x\oplus y)+s(x\odotop y)$ (where $\oplus$ is the strong
disjunction); since $s(x\oplus y)\le 1$, this gives
$s(x)+s(y)\le 1+s(x\odotop y)$, i.e.\ (S4).  For a product the four axioms hold
coordinatewise and are preserved by the convex combination
$s=\sum_i \lambda_i s_i$ with weights $\lambda_i=h_i(\Top_i)/h(\Top)\ge0$,
$\sum_i\lambda_i=1$, where $s_i$ is the factor state and the product height is the
sum of factor heights.  Hence in all listed cases $s$ is a state, proving
\eqref{it:exist-state}.

\eqref{it:exist-topf} If $s(e)=1$ then $h(e)=h(\Top)$; as $e\le\Top$, any maximal
chain from $\Bot$ to $e$ extends to one from $\Bot$ to $\Top$, and equal length
forces $e=\Top$.  The converse is $s(\Top)=1$.
\end{proof}

\begin{theorem}[$\rho_s$ is a metric, not merely a pseudometric]\label{thm:metric}
Let $s$ be a top-faithful state on a finite $\bL$ \emph{(}$s(e)=1\Rightarrow
e=\Top$\emph{)} and $\mu_n$ the uniform counting measure on $L^n$.  Then on
$\Form/{\approx}$ the logic distance $\rho_s$ is a genuine metric:
$\rho_s(\varphi,\psi)=0\iff \varphi\approx\psi$.  By
Proposition~\ref{prop:exist}\eqref{it:exist-topf} such a state exists for every
finite Heyting algebra \emph{(}including all finite distributive non-chain
lattices\emph{)}, every finite MV-algebra, and their products; on these $\rho_s$
separates points.
\end{theorem}

\begin{proof}
``$\Leftarrow$'' is the last clause of Theorem~\ref{thm:pseudometric}.  For
``$\Rightarrow$'', suppose $\varphi\not\approx\psi$.  By
Definition~\ref{def:eqv} there is a point $a\in L^m$ with
$x\eqdef\overline\varphi(a)\ne\overline\psi(a)\eqdef y$.  By
Lemma~\ref{lem:res}\eqref{it:bi-refl}, $x\fbi y\ne\Top$, hence by top-faithfulness
$s(x\fbi y)<1$, i.e.\ $1-s(\overline{(\varphi\fbi\psi)}(a))>0$.  The integrand in
\eqref{eq:rho-s} is nonnegative, and the uniform measure assigns mass
$\mu_m(\{a\})=|L|^{-m}>0$ to the singleton $\{a\}$, so
\[
\rho_s(\varphi,\psi)\ \ge\ \mu_m(\{a\})\,\bigl(1-s(x\fbi y)\bigr)\ =\
|L|^{-m}\bigl(1-s(x\fbi y)\bigr)\ >\ 0 .
\]
\end{proof}

\begin{remark}[The precise dichotomy]\label{rem:obstr-scope}
The picture requested in part (c) is the following.
\begin{itemize}
\item For \emph{every} complete residuated lattice the intrinsic, $\bL$-valued
theory of \S\ref{sec:meta} (Theorem~\ref{thm:intrinsic}) holds verbatim; it uses
only completeness and residuation, separates $\not\approx$-classes (clause
\eqref{it:E-sep}), and needs no measure or numerical structure.  This is the most
general satisfactory generalization, and its codomain is $\bL$ itself.
\item A \emph{numerical}, $[0,1]$-valued pseudometric $\rho_s$ exists whenever
$\bL$ carries a state; by Proposition~\ref{prop:exist} this includes every finite
Heyting algebra and MV-algebra and their products, in particular non-chain
lattices such as $M_2$ and $W_3\times W_4$.  For a top-faithful state $\rho_s$ is
in fact a genuine metric on $\Form/{\approx}$ (Theorem~\ref{thm:metric}).
\item The state hypothesis is sharp: replacing it by mere monotonicity destroys
the triangle inequality (Proposition~\ref{prop:obstr-mono}).  Thus the modular
law (S3) together with $\&$-subadditivity (S4) is exactly the additional
structure on $\bL$ that lets the integral/averaging construction of the classical
$[0,1]$-valued quantitative logic survive the passage to a non-linear truth-value
lattice.
\end{itemize}
When $\bL=[0,1]$ (or $W_n$) with the \L ukasiewicz operations, $s=\mathrm{id}$ is
a top-faithful state and $\mu_n$ the uniform/Lebesgue measure, recovering exactly
the classical theory of the Background paragraph.
\end{remark}

\section{Summary}\label{sec:summary}

\begin{theorem}[The lattice-valued quantitative logic framework]\label{thm:main}
Let $\bL$ be a complete residuated lattice and $\Form$ the formula algebra of
\S\ref{sec:formulas}.  Then:
\begin{enumerate}
\item\emph{(part (a))} The connectives $\wedge,\vee,\odotop,\fto$ interpreted by
the operations of $\bL$ give every $\varphi\in\Form_n$ a well-defined induced
truth function $\overline\varphi:L^n\to L$
(Definitions~\ref{def:form}--\ref{def:val}).
\item\emph{(part (b), intrinsic)} The infimum biresiduum
$E(\varphi,\psi)=\tau_\wedge(\varphi\fbi\psi)\in L$ is a well-defined
$\bL$-valued meta-degree; it satisfies reflexivity, symmetry, residuated
transitivity, separation up to $\approx$, and congruence with the connectives
(Theorem~\ref{thm:intrinsic}).  No measure on $\bL$ is required; the codomain is
$\bL$ itself.
\item\emph{(part (b), numerical)} If $\bL$ carries a state $s$
(Definition~\ref{def:state}), then $\tau_s(\varphi)\in[0,1]$ of
\eqref{eq:tau-s} is a well-defined numerical truth degree, invariant under
$\approx$ (Lemma~\ref{lem:tau-wd}); the state plays exactly the role of the
uniform/Lebesgue measure of the classical theory and the codomain is $[0,1]$.
\item\emph{(part (c))} The induced distance $\rho_s(\varphi,\psi)=1-\tau_s(\varphi\fbi\psi)$
is a pseudometric on $\Form$, descends to $\Form/{\approx}$
(Theorem~\ref{thm:pseudometric}), all connectives are non-expansive for it
(Proposition~\ref{prop:lipschitz}), and it supports a consistent theory of
divergence and consistency degrees of theories
(Proposition~\ref{prop:cons}).  All these properties reduce to the classical
$[0,1]$-valued ones when $\bL=[0,1]$ (or $W_n$) with the \L ukasiewicz operations
and $s=\mathrm{id}$, $\mu_n=$ uniform measure.
\item\emph{(sharpness and metricity)} The state hypothesis is sharp: an
arbitrary monotone collapse $\bL\to[0,1]$ that is not a state violates the
triangle inequality (Proposition~\ref{prop:obstr-mono}).  Conversely, states
exist for every finite Heyting algebra, every finite MV-algebra, and their
products (Proposition~\ref{prop:exist}), in particular for non-chain lattices, and
a top-faithful state makes $\rho_s$ a genuine point-separating metric on
$\Form/{\approx}$ (Theorem~\ref{thm:metric}).  For every complete residuated
lattice the intrinsic $\bL$-valued theory of clause (2) holds unconditionally.
\end{enumerate}
\end{theorem}

\begin{proof}
Each clause is the cited result: (1) Definitions~\ref{def:form}--\ref{def:val};
(2) Theorem~\ref{thm:intrinsic}; (3) Definition~\ref{def:tau-state} with
Lemma~\ref{lem:tau-wd}; (4) Theorems~\ref{thm:pseudometric},
Proposition~\ref{prop:lipschitz}, Proposition~\ref{prop:cons}, with the classical
reduction noted in Remark~\ref{rem:state-exist}; (5)
Proposition~\ref{prop:obstr-mono}, Proposition~\ref{prop:exist},
Theorem~\ref{thm:metric}, and Remark~\ref{rem:obstr-scope}.
\end{proof}

\end{document}
